\documentclass[11pt]{article}
\usepackage[a4paper,margin=24mm]{geometry}
\usepackage[T1]{fontenc}
\usepackage{lmodern,microtype,amsmath,amssymb,amsthm,mathtools}
\usepackage[hidelinks]{hyperref}
\setlength{\emergencystretch}{2em}
\newtheorem{theorem}{Theorem}
\newtheorem{lemma}{Lemma}
\newtheorem{proposition}{Proposition}
\newtheorem{corollary}{Corollary}
\newcommand{\R}{\mathbb R}
\newcommand{\diag}{\operatorname{diag}}
\newcommand{\supp}{\operatorname{supp}}
\newcommand{\adj}{\operatorname{adj}}
\title{A determinant certificate for canonical Bell inequalities}
\author{}
\date{23 October 2026}
\begin{document}
\maketitle
\vspace{-2em}

A canonical Bell inequality separates into eight positive event terms and
one negative state component. This structure reduces positivity of its
four-dimensional state matrix to an angular determinant. That determinant
has only eight nonconstant monomials, all with nonnegative coefficients.
Their exponents form twelve elementary balanced circuits, each with an
explicit sum-of-squares identity.

We combine these two observations into a single certificate for every
canonical valid inequality and every measurement angle. Multiplication by
the determinant of the three-dimensional positive state block clears the
state elimination, including its singular limits. The resulting certificate
uses 104 fixed functions and an explicit rational semidefinite pencil of
order 133. Its coefficients depend quartically on the inequality. This
dependence is the remaining step between the certificate and an affine
semidefinite representation of the quantum set.

\section{Completing the state quadratic without division}

Use the compact half-angle cover
\[
 a_0^2+a_1^2=b_0^2+b_1^2=1,
 \qquad \psi=(W,X,Y,Z)\in S^3.
\]
Thus $A_0=B_0=\sigma_z$,
$A_1=(a_0^2-a_1^2)\sigma_z+2a_0a_1\sigma_x$, and similarly for Bob.
The circuit identities below contain cube roots of angular monomials.
To make those roots polynomial on a fixed cover, adjoin all real solutions of
\[
 x_0^6=a_0^2,\quad x_1^6=a_1^2,\quad
 y_0^6=b_0^2,\quad y_1^6=b_1^2.                       \tag{1}
\]
This compact algebraic cover includes every root sign, angle sign and
projective endpoint. Every new variable has absolute value at most one.

For the canonical source use $d\geq0$ and eight nonnegative event weights
$(\lambda,n,m,C,A,D,\alpha,\beta)$, with
\[
 \kappa=(4\lambda,4n,4m,4C,4A,4D+2d,4\alpha+2d,4\beta+2d).
\]
The eight rows below contain no source coefficients:
\[
\begin{array}{c|rrrr}
i&R_{iW}&R_{iX}&R_{iY}&R_{iZ}\\\hline
1&b_0&b_1&0&0\\
2&0&1&0&0\\
3&a_0b_1&-a_0b_0&a_1b_1&-a_1b_0\\
4&a_1&0&-a_0&0\\
5&0&0&1&0\\
6&a_1b_0&a_1b_1&-a_0b_0&-a_0b_1\\
7&0&0&b_1&-b_0\\
8&0&a_0&0&a_1
\end{array}                                                        \tag{2}
\]
The complete canonical polynomial, including the signed core, is
\[
 F=\sum_{i=1}^8\kappa_i(R_i\psi)^2-2dZ^2
   =\psi^TP\psi,\qquad
 P=\sum_i\kappa_iR_i^TR_i-2d e_4e_4^T.                \tag{3}
\]
These are precisely the homogeneous physical-parent rows, restricted to
the normalized half-angle cover.  Put $r_i=(R_{iW},R_{iX},R_{iY})^T$ and
write $P=\bigl[\begin{smallmatrix}H&b\\b^T&c\end{smallmatrix}\bigr]$.
In particular $H=\sum_i\kappa_i r_ir_i^T$.

The adjugate of $H$ is a sum of outer products of cross products of its
rows. Completing the state quadratic therefore produces the functions
\[
 f_{ijk}=\det(r_i,r_j,r_k)R_k\psi,
 \qquad 1\leq i<j\leq8,\quad k\notin\{i,j\}.          \tag{4}
\]
The $168$ indices include zero functions and linear dependencies.
Let $E=\{-1,0,1\}^2\setminus\{0\}$ and
\[
 \mathcal N=3E\ \cup\
 \{(\pm1,\pm1),(\pm1,\pm2),(\pm2,\pm1)\}\ \cup\{0\}.
                                                               \tag{5}
\]
There are exactly 21 distinct elements.  For $\nu\in\mathcal N$ define
\[
 g_\nu=4x_0^{3-\nu_1}x_1^{3+\nu_1}
           y_0^{3-\nu_2}y_1^{3+\nu_2}Z.                \tag{6}
\]
Every exponent in (6) is an integer between zero and six.

\begin{theorem}[Fixed multiplier certificate]
For every quantum-valid canonical source with $d\geq0$ and all eight
event weights nonnegative, the polynomial $(\det H)F$, on the entire
testing cover (1), belongs to the Gram cone of the fixed indexed family
\[
                    (f_{ijk})\ \cup\ (g_\nu),          \tag{7}
\]
of size 189.  A positive-semidefinite source-only Gram matrix is specified
in (9), (15)--(16) below.  If the angular determinant has a positive
stationary zero, $g_0$ is unnecessary, giving the 188-function subfamily.
No source weight is required to be positive, and no testing opening is
required to be nonzero.
\end{theorem}

The matrix completion is the following polynomial identity, valid even
when $H$ is singular:
\[
 (\det H)F=
 \sum_{i<j}\kappa_i\kappa_j
       \left(\sum_{k\notin\{i,j\}}\kappa_k f_{ijk}\right)^2
                  +(\det P)Z^2.                      \tag{8}
\]
Indeed Cauchy--Binet gives
$\adj H=\sum_{i<j}\kappa_i\kappa_j
(r_i\times r_j)(r_i\times r_j)^T$.
Also $[H\ b]\psi=\sum_k\kappa_k r_k(R_k\psi)$, and the block determinant
identity gives
$\det P=c\det H-b^T(\adj H)b$.
Substituting these three polynomial identities proves (8); neither an
inverse nor a division has been used.  Thus the first Gram block, in the
ordering (4), is the explicit block diagonal matrix
\[
 G_{\rm state}=
 \bigoplus_{i<j}\kappa_i\kappa_j\,\kappa_{\widehat{ij}}
                       \kappa_{\widehat{ij}}^T,
 \qquad \kappa_{\widehat{ij}}=(\kappa_k)_{k\notin\{i,j\}}.
                                                               \tag{9}
\]

\section{A fixed angular Gram cone for every valid source}

It remains to express the determinant term in (8) as fixed squares.
Introduce affine variables
$u=a_1/a_0$, $v=b_1/b_0$, $U=u^2$, $V=v^2$.
Write $D_u=1+u^2$, $D_v=1+v^2$.  The affine state matrix obeys
\[
 \det P_{\rm aff}=16D_u^2D_v^2u^2v^2 L(U,V),
 \qquad L=h+\sum_{e\in E}q_e U^{e_1}V^{e_2},\quad q_e\geq0.
                                                               \tag{10}
\]
For reproducibility, set
\[
\begin{array}{ll}
t_1=d+2(\lambda+n+D),&t_2=d+2(\alpha+m+A),\\
t_3=d+2(\lambda+C+D),&t_4=d+2(\alpha+\beta+m),\\
t_5=d+2(\lambda+\beta+n),&t_6=d+2(\alpha+C+A),\\
t_7=d+2(\beta+n+m),&t_8=d+2(C+A+D).
\end{array}
\]
The coefficient indexed by $e=(i-1,j-1)$ is
\[
\begin{array}{ll}
q_{00}=4\lambda\alpha t_7t_8,&
q_{01}=4(\lambda Dt_7t_6+\alpha mt_5t_8),\\
q_{02}=4mDt_5t_6,&
q_{10}=4(\alpha At_7t_3+\lambda nt_8t_4),\\
q_{12}=4(\beta mt_1t_6+CDt_5t_2),&
q_{20}=4nAt_3t_4,\\
q_{21}=4(\beta nt_3t_2+CAt_1t_4),&
q_{22}=4\beta Ct_1t_2,
\end{array}
\]
and $h=\det P_{\rm aff}(1,1)/256-\sum_e q_e$.
Expanding the determinant of (3) gives these coefficients. Validity of
the Bell inequality implies
$L(U,V)\geq0$ for all $U,V>0$.

\begin{lemma}[The 21-monomial Laurent cone]
Every $L=h+\sum_{e\in E}q_e U^{e_1}V^{e_2}$ with $q_e\geq0$ and
$L\geq0$ on the positive quadrant is a sum of squares in the fixed span
\[
             \{U^{\nu_1/6}V^{\nu_2/6}:\nu\in\mathcal N\}.
                                                               \tag{11}
\]
The source-only algebraic coefficient prescription below also covers
zero outer coefficients and minima attained only at infinity.
\end{lemma}

\begin{proof}
If $h\geq0$, use $\sqrt h$ and $\sqrt{q_e}U^{e_1/2}V^{e_2/2}$.
Suppose $h<0$ and put $I=\{e:q_e>0\}$.  The origin belongs to
$\operatorname{conv}I$: otherwise strict separation and a ray in logarithmic
coordinates make every nonconstant term tend to zero, contradicting
$h<0$.  Let $J$ consist of the exponents on the minimal face of
$\operatorname{conv}I$ containing the origin.  An exposing ray sends
every off-face monomial to zero while keeping the face monomials fixed.
Consequently
\[
              h+\sum_{e\in J}q_eU^{e_1}V^{e_2}\geq0. \tag{12}
\]
In the span of $J$, the function
$\sum_{e\in J}q_e\exp(e\cdot y)$ is coercive and strictly convex:
the origin lies in the relative interior of $\operatorname{conv}J$, and
its Hessian is $\sum q_e\exp(e\cdot y)ee^T$, positive on that span.
It therefore has a unique minimum modulo the orthogonal complement.
Choose a positive minimizer $(U_c,V_c)$ and put
\[
 w_e=q_eU_c^{e_1}V_c^{e_2}\ (e\in J),\quad w_e=0\ (e\notin J),\qquad
 c=\sum_{e\in J}w_e,\qquad r=h+c\geq0.
\]
Then $\sum w_ee=0$.  For a rank-two face, $(U_c,V_c)$ is the unique
positive solution of the two equations
\[
 \sum_{e\in J} e_1q_e U_c^{e_1+1}V_c^{e_2+1}=0,
 \qquad
 \sum_{e\in J} e_2q_e U_c^{e_1+1}V_c^{e_2+1}=0.       \tag{13}
\]
These have bidegree at most $(2,2)$ and prescribe algebraic functions
of the source coefficients.  For a rank-one face $J=\{e,-e\}$ choose
$e$ with first nonzero component positive, set
$t=\sqrt{q_{-e}/q_e}$, and take $(U_c,V_c)=(t,1)$ if $e_1=1$,
or $(1,t)$ if $e=(0,1)$.  This gives $U_c^{e_1}V_c^{e_2}=t$.

The face selection needs only fixed rational sign cases.  Order the four
directions as $(1,0)$, $(0,1)$, $(1,1)$, $(1,-1)$, and for each direction
try $\sigma=1$ before $\sigma=-1$.  Test whether $q_e,q_{-e}>0$ and
$q_f=0$ for every $f$ with
$\sigma\det(e,f)<0$.  The first successful test gives
$J=\{e,-e\}$.  If none succeeds, $J=I$ is the rank-two case.
These are exhaustive because a proper face through the origin must
be one of the four opposite pairs in $E$.

The nonnegative balanced vector $w$ is a conic combination of exactly
twelve primitive positive circuits: four opposite pairs with weights
$(1,1)$, four triangles
$\{(s,t),(-s,0),(0,-t)\}$ with weights $(1,1,1)$, and four triangles
formed by adjacent corners and the opposite axial point, with weights
$(1,1,2)$.  Here $s,t\in\{-1,1\}$.
A deterministic prescription chooses the first supported circuit in
the fixed order in the accompanying JSON file, subtracts
$\lambda_C b^C$ with
$\lambda_C=\min_{b_e^C>0} w_e/b_e^C$, and repeats.  Each subtraction
eliminates a positive entry, so at most eight steps are needed.

Put $p_e=(U/U_c)^{e_1}(V/V_c)^{e_2}$.
An opposite pair contributes $(p_e^{1/2}-p_{-e}^{1/2})^2$.
For a triangle of weight $(1,1,1)$ use the nine squares
\[
 \frac12\sum_{h\in C}\ \sum_{\{l,m\}\subset C}
       \bigl[p_h^{1/6}(p_l^{1/3}-p_m^{1/3})\bigr]^2
                         =\sum_{e\in C}p_e-3.        \tag{14}
\]
For weights $(1,1,2)$, with $e+f+2g=0$, use
$(p_e^{1/2}-p_f^{1/2})^2+2(p_{-g}^{1/2}-p_g^{1/2})^2
=p_e+p_f+2p_g-4$.
All exponents in these roots belong to (11): in (14) they are
$(h+2l)/6$, and the possible numerators are precisely among (5).
Multiply each circuit identity by $\lambda_C$, and add the constant
$r$ and the off-face monomials.  Their sum is exactly $L$.
\end{proof}

For an entirely explicit Gram prescription, let $e_\nu$ denote the
coordinate vector indexed by $\nu\in\mathcal N$.  The twelve fixed
rational positive-semidefinite matrices are
\[
\begin{array}{ll}
A_C=(e_{3e}-e_{-3e})(e_{3e}-e_{-3e})^T,
     &C=\{e,-e\},\\[2mm]
A_C=\frac12\displaystyle\sum_{h\in C}\sum_{\{l,m\}\subset C}
 (e_{h+2l}-e_{h+2m})(e_{h+2l}-e_{h+2m})^T,
     &C\text{ has weights }(1,1,1),\\[2mm]
A_C=(e_{3e}-e_{3f})(e_{3e}-e_{3f})^T
 +2(e_{-3g}-e_{3g})(e_{-3g}-e_{3g})^T,
     &e+f+2g=0.
\end{array}                                                        \tag{15}
\]
Their entries are evaluated in \texttt{multiplier\_schema.json}.
Set $T=\diag(U_c^{-\nu_1/6}V_c^{-\nu_2/6})$.  The angular Gram matrix is
\[
 G_{\rm ang}=
 \begin{cases}
 h e_0e_0^T+\sum_{e\in E}q_e e_{3e}e_{3e}^T,&h\geq0,\\
 r e_0e_0^T+\sum_{e\notin J}q_e e_{3e}e_{3e}^T
       +T(\sum_C\lambda_CA_C)T,&h<0.
 \end{cases}                                                     \tag{16}
\]
It is positive semidefinite and depends only on the source.  Its
coefficients use finite support and sign cases, positive algebraic
roots, the unique positive solution (13), and finitely many minima.
There is no testing-dependent coefficient choice.

On $a_0a_1b_0b_1\ne0$ the homogeneous form of (10) is
\[
 \det P=16a_0^2a_1^2b_0^2b_1^2
 L(a_1^2/a_0^2,b_1^2/b_0^2).
\]
On the nonnegative-root branch, multiplication of (11) by
$4|a_0a_1b_0b_1|Z$ gives exactly (6), so
$\det P\,Z^2=g^TG_{\rm ang}g$ there.
The root-sign and endpoint statements are algebraic.  In every dyad in
(15), the two indices are congruent coordinatewise modulo two.
Consequently every nonzero entry of $G_{\rm ang}$ joins indices with
the same parities.  Each product $g_\nu g_\mu$ appearing in the quadratic
form therefore has even exponents in all four new variables.  The
quadratic identity is unchanged by any of their sign flips.
In each fixed Laurent identity
set $U=(x_1/x_0)^6$, $V=(y_1/y_0)^6$ and multiply by
$16x_0^6x_1^6y_0^6y_1^6Z^2$.  Every exponent is nonnegative, because
each index lies in $[-3,3]^2$.  This gives a polynomial identity also
at zero coordinates, and $x_i^6=a_i^2$, $y_i^6=b_i^2$ identify its
right-hand side with the homogeneous determinant.  Thus every zero
opening and projective endpoint is included without division.
Together with (8)--(9)
this proves the theorem with Gram matrix
$G_{\rm state}\oplus G_{\rm ang}$.
For a positive stationary zero of $L$, take its contact as $(U_c,V_c)$,
$J=I$, and $r=0$; the constant angular feature is then absent.

\paragraph{Independent real closed fields.}
The Laurent lemma, including the finite support split and the
prescriptions (13)--(16), is first-order over the reals after its
source-side positive roots are represented by polynomial equations and
sign conditions.
It therefore holds over any real closed source field $K_c$ extending
$\R$.  The Gram coefficients and the square roots of the nonnegative
source quantities used in their dyad factorizations belong to $K_c$. Equations
(8), (14), and the polynomial homogenization just given then hold in
$K_c\otimes_{\R}K_t$ for an independent real closed testing extension
$K_t$ of $\R$.
The resulting tensor roots use exactly the fixed testing functions (7);
no source-dependent testing radical or bounded-residue choice occurs.

\begin{corollary}[Exact quartic pullback, with both inclusions]
Among canonical tuples with $d$ and all eight event weights nonnegative,
physical Bell validity is
equivalent to membership of $(\det H)F$ in the fixed Gram cone (7).
Consequently the canonical valid-slack cone is the inverse image of a
fixed rational Gram cone under an explicit homogeneous quartic
coefficient map.
\end{corollary}

\begin{proof}
The theorem proves the forward implication.  Conversely, a Gram
representation makes $(\det H)F$ nonnegative at every test.  If $d>0$,
the signed-core part of the affine $Z=0$ block is
\[
 2d\{(uW+uvX-Y)^2+(1+u^2)v^2Y^2+(1+v^2)X^2\}.
\]
It is strictly positive for every nonzero $(W,X,Y)$ when $uv\ne0$;
hence $\det H>0$ on a dense open set of tests.  There $F\geq0$, and
continuity extends the inequality to every projective test.  If $d=0$,
the original parent (3) is already a sum of nonnegative event squares.
The Jordan block reduction for two binary measurements, proved in
Chapter~2 of the companion, reduces general tests to two-dimensional
blocks.  Local basis choices place each pair in the real plane used
here.  Since $P$ is real, its value on a complex state is the sum of
its values on the real and imaginary parts.  Thus the displayed
real-state tests establish validity on the full quantum set.
Finally $H$ and $F$ are linear in the source, so $(\det H)F$ has source
degree four.  Every testing function in (7) is a polynomial with
rational coefficients on (1), so all Gram product equations are
rational and finite.
\end{proof}


\section{The rational pencil}

All data in this section are evaluated over $\mathbb Q$ in the accompanying
JSON files.  Reduce polynomial feature products by the six monic rules
\[
\begin{gathered}
a_0^2=1-x_1^6,\quad a_1^2=x_1^6,\quad
b_0^2=1-y_1^6,\quad b_1^2=y_1^6,\\
x_0^6=1-x_1^6,\qquad y_0^6=1-y_1^6.
\end{gathered}                                                    \tag{17}
\]
Their leading monomials are pairwise coprime, so they form a monic
Gr\"obner basis for the angle-cover equations.  Every feature product
is homogeneous quadratic in the state; its vanishing on $S^3$ is
equivalent to its vanishing on every state.  No state-sphere reduction
is needed.  These normal forms introduce no spurious functional
identities: the real sign sheets over the open arcs of the two Fermat
curves are Zariski dense.  Equivalently, sign averaging separates the
four half-angle parities, and the remaining reduced polynomials vanish
only when all their coefficients vanish on the two curves.

The 189 indexed functions have 175 distinct normal-form monomials and
coefficient rank 104: 83 polynomial functions and 21 angular functions.
The exported 104-column basis and its exact rational reconstruction
matrix express every indexed function in that basis.  In the 28 state
blocks of (9), two blocks vanish, 22 have function rank four, and four
have rank six.  Congruence by their exported rational basis maps gives
the fixed structured Gram cone with block orders
\[
       \underbrace{4,\ldots,4}_{22},\quad
       \underbrace{6,\ldots,6}_{4},\quad21.             \tag{18}
\]
Its total order is $133$ and it has $535$ symmetric Gram coordinates.
The exact source blocks are pushed forward by the same congruences;
positivity and every represented polynomial are preserved.

Let $z\in\R^{535}$ be those upper-triangular Gram coordinates.  With
off-diagonal products counted twice, the evaluated product matrix is
\[
       \mathcal A\in\mathbb Z^{261\times535},\qquad
       \operatorname{rank}\mathcal A=189.              \tag{19}
\]
It has 5653 nonzero entries.  Choose the 189 pivot rows $J$ and pivot
columns $P$ recorded in the export, and let $K$ denote the other 346
columns.  Put $B=\mathcal A_{J,P}$.  All coefficient equations are
eliminated exactly by
\[
 z_K=t,\qquad z_P=B^{-1}y-B^{-1}\mathcal A_{J,K}t.       \tag{20}
\]
Inserting (20) into the block matrices (18) gives an actual rational
pencil $\mathcal L(y,t)$ of order 133, with 189 visible coefficient
variables, 346 free Gram variables, and 6079 nonzero upper-triangle
coefficient records.  The export includes the coefficient reconstruction
$\mathcal A_{:,P}B^{-1}y$, so all 261 product equations are specified,
including their 72 linear dependencies.  Conversely, every feasible
Gram tuple is represented by (20), because its $K$ coordinates are free
and $B$ is invertible.  Thus this pencil describes the structured Gram
cone exactly, in both directions.

For a canonical source $c=(d,\lambda,n,m,C,A,D,\alpha,\beta)$, let
\[
 y(c)=\operatorname{coeff}_{J}\bigl((\det H_c)F_c\bigr).
                                                               \tag{21}
\]
The evaluated map (21) has 7472 rational monomial records, all of
homogeneous source degree four.  It is obtained from (9), with the
angular target initially represented by the diagonal entries $h,q_e$;
this initial diagonal representation is an identity, not a claim that
$h$ is nonnegative.  The circuit construction supplies a positive Gram
when the source is valid.  The preceding theorem and its reverse
implication give the exact statement
\[
\boxed{\quad c\geq0:\quad
 F_c\text{ is quantum-valid}\quad\Longleftrightarrow\quad
 \exists t\in\R^{346}\;\mathcal L(y(c),t)\succeq0.\quad}           \tag{22}
\]
For a fixed source, (22) is finite semidefinite feasibility.  The input
$y(c)$ is quartic.  Treating its monomials as free lift variables would
drop their defining relations and would not preserve (22).  In
particular, the 189 coefficient variables are not the eight quantum
expectations.  Formula (22) is an exact polynomial pullback of a
spectrahedral shadow, not an affine lift of $Q_{222}$.

\section{Removing the determinant multiplier}

The source-dependent multiplier in (8) cannot simply be removed by a
fixed finite family of exact inverse or whitening factors, even on
compact regions with uniform positive reserve.  Here is a physical
example.

Take $d=1$, $\lambda=n=m=C=\alpha=\beta=1$, $D=2$, and $A=a\in[1,2]$.
For the CHSH convention $S=E_{00}-E_{01}+E_{10}+E_{11}$, the sum of the
eight losing event probabilities is $2-S/2$.  Hence this Bell slack is
\[
 f_a=9-\frac52S+4(a-1)p_A+4p_D
                    \geq9-5\sqrt2>\frac32.             \tag{23}
\]
The inequality follows directly from $S^2=4I+[A_0,A_1]\otimes[B_0,B_1]$
and $\|[A_0,A_1]\|,\|[B_0,B_1]\|\leq2$.
Thus every source in this compact family is strictly quantum valid.

At the noncommuting affine tests $u=1$ and $v\in[1/2,1]$, its $Z=0$
block is exactly
\[
 H_a(v)=\begin{pmatrix}
 8v^2+22&14v&-14\\
 14v&32v^2+18&-14v\\
 -14&-14v&8a(1+v^2)+20v^2+14
 \end{pmatrix}.                                           \tag{24}
\]
In particular $\frac{15}{4}I\preceq H_a(v)\preceq94I$ throughout this
compact source--test rectangle.  The lower bound follows from (23) and
$F=2(1+v^2)\|\psi\|^2f$; the upper bound follows from the absolute row
sums of (24).  Direct evaluation gives
\[
 (H_a(v)^{-1})_{YY}=
 \frac{1}{8(1+v^2)(a+r(v))},\qquad
 r(v)=\frac{320v^4+621v^2+126}{2(1+v^2)(64v^2+99)}.       \tag{25}
\]
The rational function $r$ is nonconstant; for example,
$r(1/2)=241/230$ and $r(1)=1067/652$.

\begin{proposition}[Rank-one inverse obstruction]
For any distinct $a_1,\ldots,a_N\in[1,2]$, the functions
$(H_{a_j}^{-1})_{YY}$ are linearly independent on $[1/2,1]$.
Consequently no finite-dimensional real space of testing functions
can provide exact inverse Gram factors $H_a^{-1}=B_aB_a^T$ for every
$a\in[1,2]$.  In particular no such space supports square invertible
exact whitening matrices $C_a^TH_aC_a=I$ for this family.
\end{proposition}

\begin{proof}
After multiplication by $8(1+v^2)$, a dependence would have the form
$\sum_jc_j/(a_j+r(v))=0$.  Multiplying denominators produces a
polynomial $Q(r(v))=0$.  Since a nonconstant rational function takes
infinitely many values on an interval, $Q=0$.  Evaluating $Q(-a_j)$
then gives $c_j\prod_{k\ne j}(a_k-a_j)=0$, so every $c_j=0$.
If all entries of all $B_a$ belonged to a fixed finite space $V$, every
entry of $B_aB_a^T$ would belong to the finite-dimensional product span
$\operatorname{span}\{vw:v,w\in V\}$, contradicting this independence.
Finally $C_a^TH_aC_a=I$ with $C_a$ square invertible implies
$H_a^{-1}=C_aC_a^T$.
\end{proof}

This proposition rules out that exact-whitening construction.  It does
not rule out other fixed Gram representations of the original slack,
and it is not an obstruction to the existence of a finite lift of the
quantum set.  In fact, the example itself has a fixed unmultiplied
certificate: with $T_0=A_0-(B_0-B_1)/\sqrt2$ and
$T_1=A_1-(B_0+B_1)/\sqrt2$, the identity
$T_0^2+T_1^2=\sqrt2(2\sqrt2I-S)$ gives
\[
 9I-\tfrac52S=(9-5\sqrt2)I+
                 \tfrac5{2\sqrt2}(T_0^2+T_1^2).
\]
Its two additional event terms are already positive squares.
The precise unfinished step is to obtain a source-independent
finite Gram cone for $F$ itself from an identity such as (8), without
dividing by the mixed source--testing function $\det H$ or introducing
source-dependent testing functions.  At some zero-weight boundary
sources the multiplier also vanishes at commuting tests; the continuous
identity proved above does not turn that zero into an unmultiplied
certificate.
\end{document}
